\chapter{1996年上海卷（理）}
\section{单选题}

\begin{problem}
已知集合 \(M = \{(x, y) | x + y = 2\}, N = \{(x, y) | x - y = 4\}\)，那么集合 \(M \cap N\) 为（\quad）

\choices
    {\(x = 3, y = -1\)}
    {\((3,-1) \)}
    {\(\{3, -1\} \)}
    {\(\{(3,-1)\} \)}

\end{problem}

\begin{problem}
在下列各区间中，函数 \(y = \sin\left(x + \frac{\pi}{4}\right) \) 的单调递增区间是（\quad）

\choices
    {\(\left[\frac{\pi}{2},\pi\right] \)}
    {\(\left[0,\frac{\pi}{4}\right] \)}
    {\([-\pi, 0]\)}
    {\(\left[\frac{\pi}{4},\frac{\pi}{2}\right] \)}

\end{problem}

\begin{problem}
如果 \(\log_{2}3 > \log_{2}3 > 0 \) ，那么 \(a\)、\(b\) 间的关系是（\quad）

\choices
    {\(0 < a < b < 1\)}
    {\(1 < a < b\)}
    {\(0 < b < a < 1\)}
    {\(1 < b < a\)}

\end{problem}

\begin{problem}
在下列命题中，真命题是（\quad）

\choices
    {若直线 \(m\)、\(n\) 都平行于平面 \(\alpha \)，则 \(m\parallel n\)}
    {设 \(\alpha-l-\beta\) 是直二面角，若直线 \(m\perp l\)，则 \(m\perp\beta\)}
    {若直线 \(m\)、\(n\) 在平面 \(\alpha\) 内的射影依次是一个点和一条直线，且 \(m\perp n\)，则 \(n\) 在 \(\alpha\) 内或 \(n\) 与 \(\alpha\) 平行}
    {设 \(m\)、\(n\) 是异面直线，若 \(m\) 与平面 \(\alpha\) 平行，则 \(n\) 与 \(\alpha\) 相交}

\end{problem}

\begin{problem}
将椭圆 \(\frac{x^2}{25} + \frac{y^2}{9} = 1\) 绕其左焦点按逆时针方向旋转 \(90^\circ\) 后所得的椭圆方程是（\quad）

\choices
    {\(\frac{(x + 4)^2}{25} +\frac{(y - 4)^2}{9} = 1\)}
    {\(\frac{(x + 4)^2}{25} +\frac{(y + 4)^2}{9} = 1\)}
    {\(\frac{(x + 4)^2}{9} +\frac{(y - 4)^2}{25} = 1\)}
    {\(\frac{(x + 4)^2}{9} +\frac{(y + 4)^2}{25} = 1\)}

\end{problem}

\begin{problem}
若函数 \(f(x)\)、\(g(x)\) 的定义域和值域都为 \(\R\)，则 \(f(x)>g(x)\)（\(x\in\R\)）成立的充要条件是（\quad）

\choices
    {有一个 \(x\in\R\)，使得 \(f(x)>g(x)\)}
    {有无穷多个 \(x\in\R\)，使得 \(f(x)>g(x)\)}
    {对 \(\R\) 中任意的 \(x\)，都有 \(f(x)>g(x)+1\)}
    {\(\R\) 中不存在 \(x\)，使得 \(f(x)\leqslant g(x)\)}

\end{problem}

\begin{problem}
若 \(\theta \in \left[0, \frac{\pi}{2}\right]\)，则椭圆 \(x^{2} + 2y^{2} - 2\sqrt{2}x\cos \theta + 4y\sin \theta = 0\) 的中心轨迹是（\quad）

\choices
    {}
    {}
    {}
    {}

\end{problem}

\begin{problem}
在下列图象中，二次函数 \(y = ax^2 + bx \) 与指数函数 \(y = \left(\frac{b}{a}\right)^x \) 的图象只可能是（\quad）

\choices
    {}
    {}
    {}
    {}

\end{problem}

\section{填空题}

\begin{problem}
方程 \(\log_{2}(9^{x}-5)=\log_{2}(3^{x}-2)+2\) 的解是 \(x=\)\fillinblank{}.

\end{problem}

\begin{problem}
函数 \(y=\frac{1}{\sqrt{\log_{\frac{1}{2}}(2-x)}} \) 的定义域是 \fillinblank{}.

\end{problem}

\begin{problem}
方程 \(\sin\frac{x}{2}=\frac{1}{3}\) 在 \([\pi,2\pi]\) 上的解是 \(x=\)\fillinblank{}.

\end{problem}

\begin{problem}
函数 \(y=x^{-2}\)（\(x<0\)）的反函数是 \(y=\)\fillinblank{}.

\end{problem}

\begin{problem}
极坐标方程分别是 \(\rho = \cos\theta \) 和 \(\rho = \sin\theta \) 的两个圆的圆心距是 \fillinblank{}.

\end{problem}

\begin{problem}
在 \((1+x)^{6}(1-x)^{4}\) 的展开式中，\(x^{3}\) 的系数是（结果用数值表示）\fillinblank{}.

\end{problem}

\begin{problem}
已知 \(O(0,0)\) 和 \(A(6,3)\) 两点，若点 \(P\) 在直线 \(OA\) 上，且 \(\frac{OP}{PA} = \frac{1}{2}\)，又 \(P\) 是线段 \(OB\) 的中点，则点 \(B\) 的坐标是 \fillinblank{}.

\end{problem}

\begin{problem}
平移坐标轴将抛物线 \(4x^{2}-8x+y+5=0\) 化为标准方程 \(x^{\prime2}=ay^{\prime}\)（\(a\neq0\)），则新坐标系的原点在原坐标系中的坐标是 \fillinblank{}.

\end{problem}

\begin{problem}
已知 \(\overrightarrow{a} + \overrightarrow{b} = 2\overrightarrow{i} - 8\overrightarrow{j}, \overrightarrow{a} - \overrightarrow{b} = -8\overrightarrow{i} + 16\overrightarrow{j} \) . 那么 \(\overrightarrow{a} \cdot \overrightarrow{b} = \) \fillinblank{}.

\end{problem}

\begin{problem}
\(\displaystyle\lim_{x\to2}\left(\frac{4}{x^{2}-4}-\frac{1}{x-2}\right)=\) \fillinblank{}.

\end{problem}

\begin{problem}
有8本互不相同的书，其中数学书3本，外文书2本，其它书3本. 若将这些书随机地排成一列放在书架上，则数学书恰好排在 \fillinblank{}

\end{problem}

\begin{problem}
有8本互不相同的书，其中数学书3本，外文书2本，其它书3本. 若将这些书排成一列放在书架上，则数学书恰好排在一起，外文书也恰好排在一起的排法共有 \fillinblank{}种(结果用数值表示).

\end{problem}

\begin{problem}
如图，在正三角形 \(ABC\) 中，\(E, F\) 依次是 \(AB, AC\) 的中点，\(AD \perp BC\)，\(EH \perp BC\)，\(FG \perp BC\)，\(D, H, G\) 为垂足，若将正三角形 \(ABC\) 绕 \(AD\) 旋转一周所得的圆锥的体积为 \(V\)，则其中由阴影部分所产生的旋转体的体积与 \(V\) 的比值是 \fillinblank{}.

\end{problem}

\begin{problem}
若 \(\i\) 是虚数单位，计算：\(\displaystyle \frac{(1+\sqrt{3}\i)^{5}}{16\sqrt{2}\left(\cos \frac{\pi}{6}-\i \sin \frac{\pi}{6}\right)}=\) \fillinblank{}.

\end{problem}

\section{解答题}

\begin{problem}
已知 \(\sin\left(\frac{\pi}{4}+\alpha\right)\sin\left(\frac{\pi}{4}-\alpha\right)=\frac{1}{6}\)，\(\alpha\in\left(\frac{\pi}{2},\pi\right)\)，求 \(\sin4\alpha\) 的值.

\end{problem}

\begin{problem}
在如图所示的直角坐标系中，一运动物体经过点 \(A(0,9)\)，其轨迹方程为 \(y=ax^{2}+c\)（\(a<0\)），\(D=(6,7)\) 为 \(x\) 轴上的给定区间. （1）为使物体落在 \(D\) 内，求 \(a\) 的取值范围；（2）若物体运动时又经过点 \(P(2,8.1)\)，问它能否落在 \(D\) 内？并说明理由.

\end{problem}

\begin{problem}
如图，在二面角 \(\alpha-l-\beta\) 中，\(A\)、\(B\in\alpha\)，\(C\)、\(D\in l\)，\(ABCD\) 为矩形. \(P\in\beta\)，\(PA\perp\alpha\)，且 \(PA=AD\). \(M\)、\(N\) 依次是 \(AB\)、\(PC\) 的中点. （1）求二面角 \(\alpha-l-\beta\) 的大小；（2）求证：\(MN\perp AB\)；（3）求异面直线 \(PA\) 与 \(MN\) 所成角的大小.

\end{problem}

\section{选作题}

\begin{problem}
设 \(z\) 是虚数，\(w=z+\frac{1}{z}\) 是实数，且 \(-1<w<2\). （1）求 \(\left|z\right|\) 的值及 \(z\) 的实部的取值范围；（2）设 \(u=\frac{1-z}{1+z}\). 求证：\(u\) 为纯虚数；（3）求 \(w-u^{2}\) 的最小值.

\end{problem}

\begin{problem}
已知 \(A(-1,2)\) 为抛物线 \(C:y=2x^{2}\) 上的点，直线 \(l_{1}\) 过点 \(A\)，且与抛物线 \(C\) 相切. 直线 \(l_{2}:x=a\)（\(a\neq-1\)）交抛物线 \(C\) 于点 \(B\)，交直线 \(l_{1}\) 于点 \(D\). （1）求直线 \(l_{1}\) 的方程；（2）设 \(\triangle ABD\) 的面积为 \(S_{1}\)，求 \(|BD|\) 及 \(S_{1}\) 的值；（3）设由抛物线 \(C\)、直线 \(l_{1}\)、\(l_{2}\) 所围成的图形的面积为 \(S_{2}\)，求证：\(S_{1}:S_{2}\) 的值为与 \(a\) 无关的常数.

\end{problem}

\begin{problem}
已知双曲线 \(S\) 的两条渐近线过坐标原点，且与以点 \(A(\sqrt{2},0)\) 为圆心、\(1\) 为半径的圆相切，双曲线 \(S\) 的一个顶点 \(A'\) 与点 \(A\) 关于直线 \(y=x\) 对称. 设直线 \(l\) 过点 \(A\)，斜率为 \(k\). （1）求双曲线 \(S\) 的方程；（2）当 \(k=1\) 时，在双曲线 \(S\) 的上丈上求点 \(B\)，使其与直线 \(l\) 的距离为 \(\sqrt{2}\)；（3）当 \(0 \leqslant k < 1\) 时，若双曲线 \(S\) 的上丈上有且只有一个点 \(B\) 到直线 \(l\) 的距离为 \(\sqrt{2}\)，求斜率 \(k\) 的值及相应的点 \(B\) 的坐标.

\end{problem}

\begin{problem}
设 \(A_{n}\) 为数列 \(\{a_{n}\}\) 前 \(n\) 项的和，\(A_{n}=\frac{3}{2}(a_{n}-1)\)（\(n\in N\)）. 数列 \(\{b_{n}\}\) 的通项公式为 \(b_{n}=4n+3\)（\(n\in N\)）. （1）求数列 \(\{a_{n}\}\) 的通项公式；（2）若 \(d\in\{a_{1}, a_{2}, a_{3}, \cdots, a_{n}, \cdots\}\cap\{b_{1}, b_{2}, b_{3}, \cdots, b_{n}, \cdots\}\)，则称 \(d\) 为数列 \(\{a_{n}\}\) 与 \(\{b_{n}\}\) 的公共项. 将数列 \(\{a_{n}\}\) 与 \(\{b_{n}\}\) 的公共项，按它们的原数列中的先后顺序排成一个新的数列 \(\{d_{n}\}\)，证明数列 \(\{d_{n}\}\) 的通项公式为 \(d_{n}=3^{2n+1}\)（\(n\in N\)）；（3）设数列 \(\{d_{n}\}\) 中的第 \(n\) 项是数列 \(\{b_{n}\}\) 中的第 \(r\) 项，\(B_{r}\) 为数列 \(\{b_{n}\}\) 前 \(r\) 项的和，\(D_{n}\) 为数列 \(\{d_{n}\}\) 前 \(n\) 项的和，\(T_{n}=B_{r}-D_{n}\). 求 \(\displaystyle\lim_{n\to\infty}\frac{T_{n}}{(a_{n})^{4}}\).
\end{problem}

